% --- Archon physics formalization source begin ---
% archon:physics
% archon:covers PhyXMiniProblems/problem_phyx_mini_0990.lean
% archon:source-report reports/phyx_mini/problem_phyx_mini_0990.source.json

\chapter{Physics problem phyx\_mini\_0990}
\label{ch:PhyXMiniProblems_problem_phyx_mini_0990}

\paragraph{Problem source.}
\#\# Physical scenario

A long solenoid with \textbackslash{}( N\_1 \textbackslash{}) windings and radius \textbackslash{}( b \textbackslash{}) surrounds a coaxial, narrower solenoid with \textbackslash{}( N\_2 \textbackslash{}) windings and radius \textbackslash{}( a < b \textbackslash{}), as shown in figure. At the far end, the outer solenoid is attached to the inner solenoid by a wire, so that current flows down the outer coil and then back through the inner coil as shown.

\#\# Question

What would be the inductannce if the sense of the inner coil were reversed?

\#\# Answer choices

- A: A: 10.0 Ω

- B: B: 6.67 Ω

- C: C: 3.33 Ω

- D: D: 13.3 Ω

\#\# Figure caption (auxiliary; use the image as primary evidence)

The image shows a solenoid consisting of a closely wound helical coil, represented in a cylindrical shape. The solenoid is illustrated with yellow windings and has an arrow labeled "I" indicating the direction of current flow into the solenoid via two leads. 

Key elements in the image include:

1. **Dimensions and Labels:**
   - The length of a single coil loop is labeled as "λ".
   - The outer radius of the solenoid is labeled as "b".
   - The inner radius is labeled as "a".
   - The number of turns (windings) is labeled as \textbackslash{}(N\_1\textbackslash{}) for the outer coil and \textbackslash{}(N\_2\textbackslash{}) for the inner winding.

2. **Text:**
   - "Leads" refers to the wires connected to the solenoid carrying current.
   - "Solenoids connected at this end" indicates a specific point on the solenoid where connections are made.

The image is likely representing the structure and function of a double-wound solenoid with emphasis on the geometry and winding configuration.

\#\# Dataset metadata

Electromagnetic Induction; Physical Model Grounding Reasoning, Spatial Relation Reasoning

\paragraph{Recorded answer/context.}
B: B: 6.67 Ω

\paragraph{Figure/image path.}
/root/proposal\_for\_physic/science-mango/phyx\_mini\_run/phyx\_data/test\_image/990.png

\paragraph{Formalization target.}
create a compiling Lean file with sorry bodies at `PhyXMiniProblems/problem\_phyx\_mini\_0990.lean`. The Lean declarations must preserve the physical quantities, dimensions or dimensional roles, figure labels, governing-law hypotheses, and final relation expressed by this problem.
Use Mathlib/Physlib names found through LeanExplore where available. If a physics API is missing, introduce faithful local abstractions rather than scalar placeholder aliases.

\begin{theorem}[Physics formalization target]
\label{thm:physics:phyx_mini_0990:target}
\lean{PhyXMiniProblems.ProblemPhyXMini0990.problem_phyx_mini_0990}
\uses{def:physics:phyx-mini-0990:phyxminiproblems-problemphyxmini0990-lengthinmeters, def:physics:phyx-mini-0990:phyxminiproblems-problemphyxmini0990-inductanceinhenries, def:physics:phyx-mini-0990:phyxminiproblems-problemphyxmini0990-permeabilityinhenriespermeter, def:physics:phyx-mini-0990:phyxminiproblems-problemphyxmini0990-coaxialseriessolenoidsetup, def:physics:phyx-mini-0990:phyxminiproblems-problemphyxmini0990-matchessuppliedcoaxialsolenoidfigure, def:physics:phyx-mini-0990:phyxminiproblems-problemphyxmini0990-hasphysicallongcoaxialgeometry, def:physics:phyx-mini-0990:phyxminiproblems-problemphyxmini0990-modelsinnerwindingreversal, def:physics:phyx-mini-0990:phyxminiproblems-problemphyxmini0990-satisfieslongcoaxialsolenoidlaws, def:physics:phyx-mini-0990:phyxminiproblems-problemphyxmini0990-satisfiesseriesmagneticenergylaws}
The assigned autoformalize agent should translate this physics problem into Lean declarations in the covered file, with theorem and lemma proof bodies written as `by sorry`.
\end{theorem}
\begin{proof}
This is an autoformalization task, not a proof task. Produce statements that can later be proved by the physics prover without weakening the physical model.
\end{proof}

% --- Archon declaration topology begin ---
\section{Declaration topology}
\begin{definition}[electric Current Dimension]
  \label{def:physics:phyx-mini-0990:phyxminiproblems-problemphyxmini0990-electriccurrentdimension}
  \lean{PhyXMiniProblems.ProblemPhyXMini0990.electricCurrentDimension}
  Electric current has dimension charge per time
\end{definition}
\begin{proof}
  This is the declared model definition of electric Current Dimension.
\end{proof}
\begin{definition}[energy Dimension]
  \label{def:physics:phyx-mini-0990:phyxminiproblems-problemphyxmini0990-energydimension}
  \lean{PhyXMiniProblems.ProblemPhyXMini0990.energyDimension}
  Energy has dimension M L² T⁻²
\end{definition}
\begin{proof}
  This is the declared model definition of energy Dimension.
\end{proof}
\begin{definition}[inductance Dimension]
  \label{def:physics:phyx-mini-0990:phyxminiproblems-problemphyxmini0990-inductancedimension}
  \lean{PhyXMiniProblems.ProblemPhyXMini0990.inductanceDimension}
  \uses{def:physics:phyx-mini-0990:phyxminiproblems-problemphyxmini0990-electriccurrentdimension, def:physics:phyx-mini-0990:phyxminiproblems-problemphyxmini0990-energydimension}
  Inductance is energy divided by current squared
\end{definition}
\begin{proof}
  This is the declared model definition of inductance Dimension.
\end{proof}
\begin{definition}[magnetic Permeability Dimension]
  \label{def:physics:phyx-mini-0990:phyxminiproblems-problemphyxmini0990-magneticpermeabilitydimension}
  \lean{PhyXMiniProblems.ProblemPhyXMini0990.magneticPermeabilityDimension}
  \uses{def:physics:phyx-mini-0990:phyxminiproblems-problemphyxmini0990-inductancedimension}
  Magnetic permeability has dimension inductance per unit length
\end{definition}
\begin{proof}
  This is the declared model definition of magnetic Permeability Dimension.
\end{proof}
\begin{definition}[electrical Resistance Dimension]
  \label{def:physics:phyx-mini-0990:phyxminiproblems-problemphyxmini0990-electricalresistancedimension}
  \lean{PhyXMiniProblems.ProblemPhyXMini0990.electricalResistanceDimension}
  \uses{def:physics:phyx-mini-0990:phyxminiproblems-problemphyxmini0990-inductancedimension}
  Resistance has dimension inductance per unit time
\end{definition}
\begin{proof}
  This is the declared model definition of electrical Resistance Dimension.
\end{proof}
\begin{definition}[Length Magnitude]
  \label{def:physics:phyx-mini-0990:phyxminiproblems-problemphyxmini0990-lengthmagnitude}
  \lean{PhyXMiniProblems.ProblemPhyXMini0990.LengthMagnitude}
  A nonnegative, unit-independent physical length
\end{definition}
\begin{proof}
  This is the declared model definition of Length Magnitude.
\end{proof}
\begin{definition}[Current Magnitude]
  \label{def:physics:phyx-mini-0990:phyxminiproblems-problemphyxmini0990-currentmagnitude}
  \lean{PhyXMiniProblems.ProblemPhyXMini0990.CurrentMagnitude}
  \uses{def:physics:phyx-mini-0990:phyxminiproblems-problemphyxmini0990-electriccurrentdimension}
  A nonnegative, unit-independent electric-current magnitude
\end{definition}
\begin{proof}
  This is the declared model definition of Current Magnitude.
\end{proof}
\begin{definition}[Energy Magnitude]
  \label{def:physics:phyx-mini-0990:phyxminiproblems-problemphyxmini0990-energymagnitude}
  \lean{PhyXMiniProblems.ProblemPhyXMini0990.EnergyMagnitude}
  \uses{def:physics:phyx-mini-0990:phyxminiproblems-problemphyxmini0990-energydimension}
  A nonnegative, unit-independent magnetic-energy magnitude
\end{definition}
\begin{proof}
  This is the declared model definition of Energy Magnitude.
\end{proof}
\begin{definition}[Inductance Magnitude]
  \label{def:physics:phyx-mini-0990:phyxminiproblems-problemphyxmini0990-inductancemagnitude}
  \lean{PhyXMiniProblems.ProblemPhyXMini0990.InductanceMagnitude}
  \uses{def:physics:phyx-mini-0990:phyxminiproblems-problemphyxmini0990-inductancedimension}
  A nonnegative, unit-independent inductance magnitude
\end{definition}
\begin{proof}
  This is the declared model definition of Inductance Magnitude.
\end{proof}
\begin{definition}[Magnetic Permeability Magnitude]
  \label{def:physics:phyx-mini-0990:phyxminiproblems-problemphyxmini0990-magneticpermeabilitymagnitude}
  \lean{PhyXMiniProblems.ProblemPhyXMini0990.MagneticPermeabilityMagnitude}
  \uses{def:physics:phyx-mini-0990:phyxminiproblems-problemphyxmini0990-magneticpermeabilitydimension}
  A nonnegative, unit-independent permeability magnitude
\end{definition}
\begin{proof}
  This is the declared model definition of Magnetic Permeability Magnitude.
\end{proof}
\begin{definition}[nonnegative SIReadout]
  \label{def:physics:phyx-mini-0990:phyxminiproblems-problemphyxmini0990-nonnegativesireadout}
  \lean{PhyXMiniProblems.ProblemPhyXMini0990.nonnegativeSIReadout}
  Read any nonnegative dimensionful quantity in coherent SI units
\end{definition}
\begin{proof}
  This is the declared model definition of nonnegative SIReadout.
\end{proof}
\begin{definition}[length In Meters]
  \label{def:physics:phyx-mini-0990:phyxminiproblems-problemphyxmini0990-lengthinmeters}
  \lean{PhyXMiniProblems.ProblemPhyXMini0990.lengthInMeters}
  \uses{def:physics:phyx-mini-0990:phyxminiproblems-problemphyxmini0990-lengthmagnitude, def:physics:phyx-mini-0990:phyxminiproblems-problemphyxmini0990-nonnegativesireadout}
  Read a length in metres
\end{definition}
\begin{proof}
  This is the declared model definition of length In Meters.
\end{proof}
\begin{definition}[current In Amperes]
  \label{def:physics:phyx-mini-0990:phyxminiproblems-problemphyxmini0990-currentinamperes}
  \lean{PhyXMiniProblems.ProblemPhyXMini0990.currentInAmperes}
  \uses{def:physics:phyx-mini-0990:phyxminiproblems-problemphyxmini0990-currentmagnitude, def:physics:phyx-mini-0990:phyxminiproblems-problemphyxmini0990-nonnegativesireadout}
  Read an electric-current magnitude in amperes
\end{definition}
\begin{proof}
  This is the declared model definition of current In Amperes.
\end{proof}
\begin{definition}[energy In Joules]
  \label{def:physics:phyx-mini-0990:phyxminiproblems-problemphyxmini0990-energyinjoules}
  \lean{PhyXMiniProblems.ProblemPhyXMini0990.energyInJoules}
  \uses{def:physics:phyx-mini-0990:phyxminiproblems-problemphyxmini0990-energymagnitude, def:physics:phyx-mini-0990:phyxminiproblems-problemphyxmini0990-nonnegativesireadout}
  Read magnetic energy in joules
\end{definition}
\begin{proof}
  This is the declared model definition of energy In Joules.
\end{proof}
\begin{definition}[inductance In Henries]
  \label{def:physics:phyx-mini-0990:phyxminiproblems-problemphyxmini0990-inductanceinhenries}
  \lean{PhyXMiniProblems.ProblemPhyXMini0990.inductanceInHenries}
  \uses{def:physics:phyx-mini-0990:phyxminiproblems-problemphyxmini0990-inductancemagnitude, def:physics:phyx-mini-0990:phyxminiproblems-problemphyxmini0990-nonnegativesireadout}
  Read an inductance in henries
\end{definition}
\begin{proof}
  This is the declared model definition of inductance In Henries.
\end{proof}
\begin{definition}[permeability In Henries Per Meter]
  \label{def:physics:phyx-mini-0990:phyxminiproblems-problemphyxmini0990-permeabilityinhenriespermeter}
  \lean{PhyXMiniProblems.ProblemPhyXMini0990.permeabilityInHenriesPerMeter}
  \uses{def:physics:phyx-mini-0990:phyxminiproblems-problemphyxmini0990-magneticpermeabilitymagnitude, def:physics:phyx-mini-0990:phyxminiproblems-problemphyxmini0990-nonnegativesireadout}
  Read a magnetic permeability in henries per metre
\end{definition}
\begin{proof}
  This is the declared model definition of permeability In Henries Per Meter.
\end{proof}
\begin{definition}[Figure Object]
  \label{def:physics:phyx-mini-0990:phyxminiproblems-problemphyxmini0990-figureobject}
  \lean{PhyXMiniProblems.ProblemPhyXMini0990.FigureObject}
  Physical objects visible in the primary raster
\end{definition}
\begin{proof}
  This is the declared model definition of Figure Object.
\end{proof}
\begin{definition}[Figure Label]
  \label{def:physics:phyx-mini-0990:phyxminiproblems-problemphyxmini0990-figurelabel}
  \lean{PhyXMiniProblems.ProblemPhyXMini0990.FigureLabel}
  Textual or symbolic annotations visible in 990.png
\end{definition}
\begin{proof}
  This is the declared model definition of Figure Label.
\end{proof}
\begin{definition}[Solenoid End]
  \label{def:physics:phyx-mini-0990:phyxminiproblems-problemphyxmini0990-solenoidend}
  \lean{PhyXMiniProblems.ProblemPhyXMini0990.SolenoidEnd}
  The two axial ends of the coaxial pair in the raster
\end{definition}
\begin{proof}
  This is the declared model definition of Solenoid End.
\end{proof}
\begin{definition}[Axial Current Direction]
  \label{def:physics:phyx-mini-0990:phyxminiproblems-problemphyxmini0990-axialcurrentdirection}
  \lean{PhyXMiniProblems.ProblemPhyXMini0990.AxialCurrentDirection}
  Directed travel along the common solenoid axis
\end{definition}
\begin{proof}
  This is the declared model definition of Axial Current Direction.
\end{proof}
\begin{definition}[Series Current Path]
  \label{def:physics:phyx-mini-0990:phyxminiproblems-problemphyxmini0990-seriescurrentpath}
  \lean{PhyXMiniProblems.ProblemPhyXMini0990.SeriesCurrentPath}
  The series-current route determined by the right-end connection
\end{definition}
\begin{proof}
  This is the declared model definition of Series Current Path.
\end{proof}
\begin{definition}[Winding Configuration]
  \label{def:physics:phyx-mini-0990:phyxminiproblems-problemphyxmini0990-windingconfiguration}
  \lean{PhyXMiniProblems.ProblemPhyXMini0990.WindingConfiguration}
  The winding state relevant to the question
\end{definition}
\begin{proof}
  This is the declared model definition of Winding Configuration.
\end{proof}
\begin{definition}[Magnetic Coupling Sense]
  \label{def:physics:phyx-mini-0990:phyxminiproblems-problemphyxmini0990-magneticcouplingsense}
  \lean{PhyXMiniProblems.ProblemPhyXMini0990.MagneticCouplingSense}
  Whether the two solenoid fields oppose or reinforce for the series current
\end{definition}
\begin{proof}
  This is the declared model definition of Magnetic Coupling Sense.
\end{proof}
\begin{definition}[Coaxial Solenoid Figure]
  \label{def:physics:phyx-mini-0990:phyxminiproblems-problemphyxmini0990-coaxialsolenoidfigure}
  \lean{PhyXMiniProblems.ProblemPhyXMini0990.CoaxialSolenoidFigure}
  \uses{def:physics:phyx-mini-0990:phyxminiproblems-problemphyxmini0990-figureobject, def:physics:phyx-mini-0990:phyxminiproblems-problemphyxmini0990-figurelabel, def:physics:phyx-mini-0990:phyxminiproblems-problemphyxmini0990-solenoidend, def:physics:phyx-mini-0990:phyxminiproblems-problemphyxmini0990-axialcurrentdirection, def:physics:phyx-mini-0990:phyxminiproblems-problemphyxmini0990-seriescurrentpath}
  Literal presentation and spatial relations transcribed from 990.png
\end{definition}
\begin{proof}
  This is the declared model definition of Coaxial Solenoid Figure.
\end{proof}
\begin{definition}[Coaxial Series Solenoid Setup]
  \label{def:physics:phyx-mini-0990:phyxminiproblems-problemphyxmini0990-coaxialseriessolenoidsetup}
  \lean{PhyXMiniProblems.ProblemPhyXMini0990.CoaxialSeriesSolenoidSetup}
  \uses{def:physics:phyx-mini-0990:phyxminiproblems-problemphyxmini0990-lengthmagnitude, def:physics:phyx-mini-0990:phyxminiproblems-problemphyxmini0990-currentmagnitude, def:physics:phyx-mini-0990:phyxminiproblems-problemphyxmini0990-energymagnitude, def:physics:phyx-mini-0990:phyxminiproblems-problemphyxmini0990-inductancemagnitude, def:physics:phyx-mini-0990:phyxminiproblems-problemphyxmini0990-magneticpermeabilitymagnitude, def:physics:phyx-mini-0990:phyxminiproblems-problemphyxmini0990-windingconfiguration, def:physics:phyx-mini-0990:phyxminiproblems-problemphyxmini0990-magneticcouplingsense, def:physics:phyx-mini-0990:phyxminiproblems-problemphyxmini0990-coaxialsolenoidfigure}
  Independent geometry, material properties, inductances, and energy observables. In particular, the reversed equivalent inductance is an independent field and is not defined from the desired closed formula
\end{definition}
\begin{proof}
  This is the declared model definition of Coaxial Series Solenoid Setup.
\end{proof}
\begin{definition}[Matches Supplied Coaxial Solenoid Figure]
  \label{def:physics:phyx-mini-0990:phyxminiproblems-problemphyxmini0990-matchessuppliedcoaxialsolenoidfigure}
  \lean{PhyXMiniProblems.ProblemPhyXMini0990.MatchesSuppliedCoaxialSolenoidFigure}
  \uses{def:physics:phyx-mini-0990:phyxminiproblems-problemphyxmini0990-coaxialseriessolenoidsetup}
  The labels, nesting, connection, and current arrows read from the image
\end{definition}
\begin{proof}
  This is the declared model definition of Matches Supplied Coaxial Solenoid Figure.
\end{proof}
\begin{definition}[Has Physical Long Coaxial Geometry]
  \label{def:physics:phyx-mini-0990:phyxminiproblems-problemphyxmini0990-hasphysicallongcoaxialgeometry}
  \lean{PhyXMiniProblems.ProblemPhyXMini0990.HasPhysicalLongCoaxialGeometry}
  \uses{def:physics:phyx-mini-0990:phyxminiproblems-problemphyxmini0990-lengthinmeters, def:physics:phyx-mini-0990:phyxminiproblems-problemphyxmini0990-permeabilityinhenriespermeter, def:physics:phyx-mini-0990:phyxminiproblems-problemphyxmini0990-coaxialseriessolenoidsetup}
  Positivity and the strict radius relation 0 < a < b
\end{definition}
\begin{proof}
  This is the declared model definition of Has Physical Long Coaxial Geometry.
\end{proof}
\begin{definition}[Models Inner Winding Reversal]
  \label{def:physics:phyx-mini-0990:phyxminiproblems-problemphyxmini0990-modelsinnerwindingreversal}
  \lean{PhyXMiniProblems.ProblemPhyXMini0990.ModelsInnerWindingReversal}
  \uses{def:physics:phyx-mini-0990:phyxminiproblems-problemphyxmini0990-coaxialseriessolenoidsetup}
  Reversing only the inner winding changes an opposing series-field arrangement into an aiding one; it does not change the geometry, turn counts, or material. This is qualitative winding data, not an inductance-value premise
\end{definition}
\begin{proof}
  This is the declared model definition of Models Inner Winding Reversal.
\end{proof}
\begin{definition}[Satisfies Long Coaxial Solenoid Laws]
  \label{def:physics:phyx-mini-0990:phyxminiproblems-problemphyxmini0990-satisfieslongcoaxialsolenoidlaws}
  \lean{PhyXMiniProblems.ProblemPhyXMini0990.SatisfiesLongCoaxialSolenoidLaws}
  \uses{def:physics:phyx-mini-0990:phyxminiproblems-problemphyxmini0990-lengthinmeters, def:physics:phyx-mini-0990:phyxminiproblems-problemphyxmini0990-inductanceinhenries, def:physics:phyx-mini-0990:phyxminiproblems-problemphyxmini0990-permeabilityinhenriespermeter, def:physics:phyx-mini-0990:phyxminiproblems-problemphyxmini0990-coaxialseriessolenoidsetup}
  The standard long-solenoid self and mutual inductances. The mutual term uses the area π a² common to both coaxial fields, not the outer area π b². None of these fields mentions the equivalent inductance after reversal
\end{definition}
\begin{proof}
  This is the declared model definition of Satisfies Long Coaxial Solenoid Laws.
\end{proof}
\begin{definition}[Satisfies Series Magnetic Energy Laws]
  \label{def:physics:phyx-mini-0990:phyxminiproblems-problemphyxmini0990-satisfiesseriesmagneticenergylaws}
  \lean{PhyXMiniProblems.ProblemPhyXMini0990.SatisfiesSeriesMagneticEnergyLaws}
  \uses{def:physics:phyx-mini-0990:phyxminiproblems-problemphyxmini0990-currentmagnitude, def:physics:phyx-mini-0990:phyxminiproblems-problemphyxmini0990-currentinamperes, def:physics:phyx-mini-0990:phyxminiproblems-problemphyxmini0990-energyinjoules, def:physics:phyx-mini-0990:phyxminiproblems-problemphyxmini0990-inductanceinhenries, def:physics:phyx-mini-0990:phyxminiproblems-problemphyxmini0990-coaxialseriessolenoidsetup}
  Magnetic energy is both L\_eq I² / 2 and the sum of the two self-energy terms plus the aiding mutual-energy term after reversal. Stating the laws for every current avoids defining the requested inductance by the target formula
\end{definition}
\begin{proof}
  This is the declared model definition of Satisfies Series Magnetic Energy Laws.
\end{proof}
\begin{definition}[Answer Label]
  \label{def:physics:phyx-mini-0990:phyxminiproblems-problemphyxmini0990-answerlabel}
  \lean{PhyXMiniProblems.ProblemPhyXMini0990.AnswerLabel}
  Labels printed beside the four dataset choices
\end{definition}
\begin{proof}
  This is the declared model definition of Answer Label.
\end{proof}
\begin{definition}[Printed Electrical Unit]
  \label{def:physics:phyx-mini-0990:phyxminiproblems-problemphyxmini0990-printedelectricalunit}
  \lean{PhyXMiniProblems.ProblemPhyXMini0990.PrintedElectricalUnit}
  The source explicitly prints the resistance unit ohm after every value
\end{definition}
\begin{proof}
  This is the declared model definition of Printed Electrical Unit.
\end{proof}
\begin{definition}[Listed Answer]
  \label{def:physics:phyx-mini-0990:phyxminiproblems-problemphyxmini0990-listedanswer}
  \lean{PhyXMiniProblems.ProblemPhyXMini0990.ListedAnswer}
  \uses{def:physics:phyx-mini-0990:phyxminiproblems-problemphyxmini0990-answerlabel, def:physics:phyx-mini-0990:phyxminiproblems-problemphyxmini0990-printedelectricalunit}
  A literal answer row, kept separate from dimensionful physical quantities
\end{definition}
\begin{proof}
  This is the declared model definition of Listed Answer.
\end{proof}
\begin{definition}[listed Answer]
  \label{def:physics:phyx-mini-0990:phyxminiproblems-problemphyxmini0990-listedanswer-value}
  \lean{PhyXMiniProblems.ProblemPhyXMini0990.listedAnswer}
  \uses{def:physics:phyx-mini-0990:phyxminiproblems-problemphyxmini0990-answerlabel, def:physics:phyx-mini-0990:phyxminiproblems-problemphyxmini0990-listedanswer}
  The four rows exactly as numerically represented in the source
\end{definition}
\begin{proof}
  This is the declared model definition of listed Answer.
\end{proof}
\begin{definition}[recorded Dataset Answer]
  \label{def:physics:phyx-mini-0990:phyxminiproblems-problemphyxmini0990-recordeddatasetanswer}
  \lean{PhyXMiniProblems.ProblemPhyXMini0990.recordedDatasetAnswer}
  \uses{def:physics:phyx-mini-0990:phyxminiproblems-problemphyxmini0990-listedanswer, def:physics:phyx-mini-0990:phyxminiproblems-problemphyxmini0990-listedanswer-value}
  The answer key records row B; this is metadata, not an inductance claim
\end{definition}
\begin{proof}
  This is the declared model definition of recorded Dataset Answer.
\end{proof}
\begin{lemma}[reversed series inductance eq self add self add twice mutual]
  \label{lem:physics:phyx-mini-0990:phyxminiproblems-problemphyxmini0990-reversed-series-inductance-eq-self-add-self-add-twice-mutual}
  \lean{PhyXMiniProblems.ProblemPhyXMini0990.reversed_series_inductance_eq_self_add_self_add_twice_mutual}
  \uses{def:physics:phyx-mini-0990:phyxminiproblems-problemphyxmini0990-inductanceinhenries, def:physics:phyx-mini-0990:phyxminiproblems-problemphyxmini0990-coaxialseriessolenoidsetup, def:physics:phyx-mini-0990:phyxminiproblems-problemphyxmini0990-hasphysicallongcoaxialgeometry, def:physics:phyx-mini-0990:phyxminiproblems-problemphyxmini0990-modelsinnerwindingreversal, def:physics:phyx-mini-0990:phyxminiproblems-problemphyxmini0990-satisfiesseriesmagneticenergylaws}
  Aiding magnetic energy gives the standard series formula
\end{lemma}
\begin{proof}
  The conclusion follows from the governing relations and figure readouts named in the statement of reversed series inductance eq self add self add twice mutual.
\end{proof}
% --- Archon declaration topology end ---
% --- Archon physics formalization source end ---
